\documentclass[12pt]{article}
\usepackage[utf8]{inputenc}
\usepackage[T1]{fontenc}
\usepackage{hyperref}
\title{Surbakhal Language Resources for K-12 Curriculum Development}
\author{Compiled by Claude Sonnet 5.5 (claude-sonnet-5-5) for Andr\'as Kornai}
\date{}
\begin{document}
\maketitle

\section{Introduction}
Surbakhal (also Sourbakhal; ISO 639-3 sbj, Glottocode surb1238) is a Maban language of eastern Chad. Its speakers live in the Ouaddai Region, Assoungha department, Adr\'e subprefecture, between Hadjer Hadid and Alacha. Speaker figures vary widely: 7,900 (2000, Wikipedia), 17,500 (Global Recordings Network) and about 25,000 for the people group (Joshua Project). Joshua Project lists the language as unwritten, with no audio Bible or Jesus Film. Searches found no school materials.

\section{Primary Educational Resources}
No primers, textbooks or school programmes were found in search results.

\section{Linguistic and Reference Materials}
\begin{itemize}
\item Surbakhal language (Wikipedia): \url{https://en.wikipedia.org/wiki/Surbakhal_language}
\item Maban languages (Wikipedia): \url{https://en.wikipedia.org/wiki/Maban_languages}
\item Maass, Grant and Prinz (1995), a sociolinguistic survey of the Sourbakhal of Ouaddai (SIL Tchad), as listed on Glottolog: \url{https://glottolog.org/resource/reference/id/111611}
\end{itemize}

\section{Religious and Community Resources}
\begin{itemize}
\item Joshua Project people-group profile: \url{https://joshuaproject.net/people_groups/15129}
\item Joshua Project language page: \url{https://joshuaproject.net/languages/sbj}
\end{itemize}

\section{Digital and Audio Resources}
\begin{itemize}
\item Global Recordings Network page for sbj, which reports no recordings currently available: \url{https://globalrecordings.net/en/language/sbj}
\end{itemize}

\section{Curriculum Development Context}
The language is unwritten and not known to be taught in schools. Masalit, a related Maban language with a larger speaker population, may be a useful comparison for orthography design (Omniglot: \url{https://omniglot.com/writing/masalit.htm}). I found no source directly comparing the two languages. The 1995 survey would be the first document to consult on language vitality and attitudes.
\end{document}